\SourceLeaf{236}%
\noindent oder wenn man
\[
\cot(431)=\alpha,\qquad\cot(234)=\beta
\]
setzt:
\[
\alpha^{2}\cot(341)^{2}-\beta^{2}\left(\frac{\operatorname{tg}i.24}{i}\right)^{2}=1.
\]
Das Tetraeder $1234$, dessen Volumen $\Delta$ heisse, möge nunmehr dadurch einen unendlich kleinen Volumenzuwachs $1241'2'4'=\partial\Delta$ erfahren, dass man die Kante $31$ um das unendlich kleine Stück $11'=d(13)$ verlängert und durch $1'$ senkrecht zu $31'$ die Ebene $1'2'4'$ legt, die die Kanten $32$ und $34$ beziehungsweise in $2'$ und $4'$ schneidet. Die Winkel an der Ecke $3$, also auch die Grössen $\alpha$ und $\beta$, bleiben hierbei unverändert. Der Winkel $(341)$ geht in den Winkel
\[
(34'1')=(341)+d(341)
\]
über, und zwar wird, wie die Betrachtung des Vierecks $11'4'4$ mit der unendlich kleinen Grundlinie $11'=d(13)$ und rechten Winkeln bei $1$ und $1'$ sofort erkennen lässt:
\[
d(341)=-\frac{\sin(i.14)}{i}.d(13).
\]
Der Volumenzuwachs $1241'2'4'=\partial\Delta$ wird seitlich von den Dreiecken $124$ und $1'2'4'$ begrenzt, deren Ebenen beide auf $11'$ senkrecht stehen, mithin ist (vergl. S.~233 dieses Bandes):
\[
\partial\Delta=\tfrac{1}{2}d(13).(24).\frac{\sin(i.14)}{i},
\]
folglich
\[
\partial\Delta=-\tfrac{1}{2}(24)d(341),
\]
und das ist, abgesehen von dem Factor $\tfrac{1}{2}$, die \textsc{Gauss}'sche Formel. Setzt man also
\[
(24)=x,\qquad(341)=\varphi,
\]
so ergibt sich für das Volumen des Tetraeders $1234$ der Ausdruck
\[
\Delta=-\tfrac{1}{2}\int_{0}^{x}x\frac{d\varphi}{dx}\,dx;
\]
dabei hat man $\dfrac{d\varphi}{dx}$ vermöge der Gleichung
\[
\alpha^{2}\cos\varphi^{2}-\beta^{2}\left(\frac{\operatorname{tg}ix}{i}\right)^{2}=1
\]
zu berechnen.
\begin{flushright}\textsc{Stäckel.}\end{flushright}
\par\endgroup
